\documentclass[11pt]{article}
\usepackage{lmodern}
\usepackage[margin=1in]{geometry}
\usepackage{amsmath,booktabs}
\title{Staffing and prices after New Jersey's 1992 minimum wage increase}
\author{}
\date{}
\begin{document}
\maketitle
\begin{abstract}
On April 1, 1992, New Jersey raised its hourly minimum wage from
\$4.25 to \$5.05. Eastern Pennsylvania restaurants faced the same
federal floor of \$4.25 as before. Across 384 incumbent fast-food
restaurants with paired staffing data, New Jersey gained 0.47
full-time equivalents per store while Pennsylvania lost 2.28. The
difference of 2.75 came chiefly from full-time nonmanagement staff.
Starting pay rose by a relative \$0.50 per hour and the price of a
three-item basket by \$0.10. Weekday opening hours in New Jersey were
almost unchanged. The estimates describe short-run adjustments at
existing locations. With one pre-policy observation and one state on
each side, the comparison cannot distinguish the policy from a
differential state employment trend.
\end{abstract}

\section{Restaurants, timing, and measurement}
The unit is a location operating at the first survey, with each
location weighted equally. Managers or assistant managers at Burger King,
KFC, Roy Rogers, and Wendy's reported staffing and operating conditions.
The two interview windows were February 15--March 4 and November
5--December 31, 1992. The increase was 18.8 percent. It was scheduled
before the first interviews, which explicitly asked New Jersey managers
whether they had heard of it. The baseline can therefore already
reflect adjustments in anticipation.

The supplied file has 331 New Jersey and 79 eastern Pennsylvania
locations. In chain order as listed above, the counts are
136, 68, 82, and 45 in New Jersey and 35, 12, 17, and 15 in Pennsylvania.
Each row contains both waves. Sheet 407 occurs once in each state,
at different chains. State and sheet jointly identify all 410 rows.
One New Jersey follow-up (sheet 178) has a recorded year of 1991.
I use the wave designation in the main analysis and check its exclusion.

The primary outcome is staffing in full-time equivalents (FTE), the
sum of full-time nonmanagers, managers and assistants, and half the
part-time nonmanager count. The survey asks for these groups separately.
The half-weight approximates work hours, which the data do not record.
Headcount gives each employee unit weight. Both measures count positions
at a location and cannot track individual workers.

I assign zero post-period employment to the six permanent closures
(five in New Jersey and one in Pennsylvania). Both states include them
in the incumbent cohort. Four New Jersey stores closed for renovations,
highway construction, or a mall fire, and one refused the follow-up.
Temporary closure leaves payroll status uncertain, so these five stores
contribute no post-period outcome in the main comparison. The other
321 New Jersey and 78 Pennsylvania restaurants answered. A dot in the
file means missing. Each outcome requires all its components at both
waves, with no imputation for item nonresponse. By state, usable
baseline/follow-up FTE counts are 321/319 and 77/77, yielding 309
and 75 complete pairs. Wage and operating outcomes require a completed
second interview and describe reporting restaurants still in operation.

\section{Design and interpretation}
The state sets the policy, and initial location assigns exposure.
Let $D_i$ denote a New Jersey store and $Y_{it}(d)$ its outcome at
wave $t$ under policy regime $d$. For the observed pairs $S$ of a
given outcome I estimate
\[
 \widehat\delta=\overline{(Y_1-Y_0)}_{D=1,S}
               -\overline{(Y_1-Y_0)}_{D=0,S}.
\]
This is the interaction coefficient in a two-wave regression with
restaurant and wave fixed effects. Differencing removes fixed
location attributes, and the Pennsylvania change measures the common
time effect. There are no post-policy control variables.

The intended causal quantity for employment is the average policy effect
on post-period staffing at the New Jersey incumbents represented by the
pairs. With a stable comparison sample, no anticipation, and no spillover
to Pennsylvania, adding and subtracting New Jersey's untreated change gives
\[
 \delta=E[Y_1(1)-Y_1(0)\mid D=1,S]+
 \underbrace{E[\Delta Y(0)\mid D=1,S]-
 E[\Delta Y(0)\mid D=0,S]}_{B}.
\]
Parallel untreated trends set $B=0$. A relative untreated New Jersey gain
of 2.75 FTE per store would explain the entire estimate. Local demand,
different winter trading patterns, or an unusually strong Pennsylvania
baseline could produce such a gap. Cross-border customer or worker
movement could also change the control outcome. Only one pre-policy
wave is available. Neither pre-trends nor the absence of spillovers
can be tested here. Extending the estimate to all incumbents requires
an assumption about missing changes. For prices and benefits,
policy-dependent survival or reporting can also change which stores have
two meaningful measurements.

\section{Staffing and pay}
\begin{table}[ht]
\centering
\caption{Paired changes per incumbent restaurant}
\small
\begin{tabular}{lrrrrrr}
\toprule
& \multicolumn{2}{c}{Paired stores} & \multicolumn{2}{c}{Mean change}
& & \\
Outcome & NJ & PA & NJ & PA & Difference & SE\\
\midrule
FTE &309&75&0.47&$-2.28$&2.75&1.34\\
Headcount &309&75&0.19&$-2.23$&2.41&1.49\\
FTE excluding managers &315&76&0.45&$-2.49$&2.94&1.32\\
Starting wage (\$/hour)&301&68&0.47&$-0.03$&0.50&0.05\\
\bottomrule
\end{tabular}
\end{table}

Mean FTE rose from 20.43 to 20.90 in New Jersey and fell from
23.38 to 21.10 in Pennsylvania. Most of the contrast comes from the
comparison state's decline. The relative gain equals about 13.5 percent
of initial New Jersey staffing. Within the identical 384 pairs, the
relative changes were $+3.16$ full-time nonmanagers, $-0.68$ part-time
nonmanagers, and $-0.07$ managers. Half-weighting the second component
recovers the FTE difference. Thus shifting the assumed part-time weight
anywhere between one-half and one preserves a positive estimate.
Changes in actual hours within employment categories remain unmeasured.

Standard errors (SE) in both tables are
$\sqrt{s^2_{\Delta,NJ}/n_{NJ}+s^2_{\Delta,PA}/n_{PA}}$, with
$s^2$ the sample variance of paired changes. They summarize precision
under independent sampling of restaurants conditional on the two state
environments. The corresponding normal-reference 95 percent interval for
FTE is $[0.12,5.38]$. Restaurants may share chain, local, or state
shocks. There are only two policy clusters, one treated. The interval
does not quantify uncertainty over state-level counterfactual trends,
and a policy-level significance claim would be unwarranted.

New Jersey starting pay rose from \$4.61 to \$5.08 versus \$4.65 to
\$4.62 in Pennsylvania (levels rounded independently). The bite of the
law differed across stores. Among New Jersey stores reporting baseline
pay, 101 started at \$4.25, 185 strictly between the two floors, and 28
at or above \$5.05. In that order, mean paired starting-pay changes were
\$0.81, \$0.38, and $-$\$0.16 (95, 179, and 27 pairs). Their FTE changes
were $+1.20$, $+0.41$, and $-4.85$ (94, 174, and 25 pairs).
Stores facing the largest required increases therefore show no average
staffing contraction. Initial pay reflects local conditions, however,
and grouping on a noisy baseline wage can itself generate convergence.
High-wage stores are few and can face competitive wage responses even
without a statutory shortfall. This comparison supplies a check on
exposure, not an independent experiment.

\section{Operating adjustments}
\begin{table}[ht]
\centering
\caption{Changes among reporting restaurants operating at follow-up}
\small
\begin{tabular}{lrrrrrr}
\toprule
& \multicolumn{2}{c}{Paired stores} & \multicolumn{2}{c}{Mean change} & & \\
Outcome & NJ & PA & NJ & PA & Difference & SE\\
\midrule
Three-item price (\$)&287&68&0.067&$-0.034$&0.101&0.044\\
Weekday hours open &320&78&0.001&0.109&$-0.108$&0.103\\
Registers open at 11 a.m.&298&74&$-0.040$&$-0.216$&0.176&0.102\\
Installed registers &306&76&$-0.042$&0.132&$-0.174$&0.108\\
Free meal offered (pp)&321&78&4.36&1.28&3.08&4.68\\
Discounted meal offered (pp)&321&78&$-8.72$&$-6.41$&$-2.31$&4.02\\
First raise (\$/hour)&241&55&$-0.010$&$-0.017$&0.008&0.019\\
\bottomrule
\end{tabular}
\end{table}

The basket sums the tax-inclusive price of one medium soda, small
fries, and regular hamburger (two pieces of chicken or equivalent at
KFC). Its New Jersey baseline was \$3.36. The relative increase of ten
cents is about 3.0 percent of that price. Within-chain price contrasts
are positive for all four chains. Weighting those contrasts by their
shares of New Jersey price pairs gives \$0.103. Higher consumer prices
are consistent with some adjustment to payroll costs. Starting wages
do not give the wage bill, and sales volumes, input costs, and net
prices before tax are unavailable for a pass-through calculation.

Weekday hours barely moved in New Jersey. The relative increase in
registers open at 11 a.m. reflects a larger Pennsylvania reduction.
Neither that snapshot nor the installed register count measures sales
or service quality. Meal indicators overlap where both benefits are
offered (pp denotes percentage points). Their mixed changes and the
small first-raise contrast give little evidence of a broad offset in
these measured benefits. First-raise amounts have fewer respondents.
The paired first-raise delay contrast is $+2.51$ in recorded units
(SE 2.16), but the forms ask for weeks and the codebook says months,
so its duration is ambiguous. The baseline recruitment field records a
cash bounty while follow-up records any special program. Their difference
would mix a change in definition with a change in practice.

\section{Robustness, mechanisms, and scope}
Excluding permanent closures gives a relative FTE change of 2.72.
Counting all four temporary closures as having zero staff gives 2.51.
Removing the anomalously dated follow-up gives 2.79, and excluding
New Jersey shore stores gives 2.97. Permanent closure rates themselves
were 5/331 and 1/79. These few events give limited information about
a policy effect on exit.

Comparing stores within the same chain and averaging contrasts with New
Jersey FTE-pair shares gives 2.50. The chain-specific contrasts are
4.39 (Burger King), $-1.59$ (KFC), 2.49 (Roy Rogers), and 3.42
(Wendy's). The aggregate pattern is not shared by every chain.
Keeping New Jersey fixed and using only the Pennsylvania northeast
Philadelphia suburbs (34 pairs) gives 4.33, versus 1.44 using the
Easton-area group (41 pairs). The sign survives these choices, but
the size depends materially on which local employment trend represents
New Jersey's counterfactual.

A competitive labor-demand response to higher wages can include reduced
hours per worker, substitution toward other inputs, or lower hiring at
new locations. None requires an immediate decline in measured incumbent
headcount if demand is also growing. Alternatively, when recruitment or
retention improves with pay, a higher wage floor can sustain or increase
staffing even as starting pay rises. The surveys lack linked worker flows,
vacancies, productivity, profits, and transaction volumes to discriminate
among these explanations. The observed price increase and staffing mix
are adjustment margins, not a decomposition of the policy's costs.

\paragraph{Conclusion.}
By late 1992 these New Jersey incumbents had raised starting wages and
menu prices relative to their Pennsylvania counterparts, with modest
absolute staffing growth and nearly unchanged weekday hours. A causal
reading of the larger relative staffing gain rests primarily on the
untested Pennsylvania counterfactual. New restaurants never enter this
baseline cohort. Estimating entry requires dated establishment openings
and a consistent geographic sampling frame. Measuring net industry or
statewide employment requires staffing for entrants and exits as well
as representative coverage beyond these chains, for example through an
establishment payroll panel. Those estimates remain unfinished because
the required data were not supplied.

\paragraph{Sources and replication.}
The institutional amounts, geography, and field windows come from
\texttt{public-task.json}. The only observations are
\texttt{input/observations.dat}. The supplied codebook gives fixed-width
columns, status codes, and units, and the two blank questionnaires give
the wording and item definitions. From the task directory, running
\texttt{python3 analysis.py} with standard-library Python regenerates
\texttt{results.json}, including unrounded estimates, denominators,
conditional precision measures, and sensitivity comparisons.
\end{document}
